% Standalone fragment: requires amsmath, amssymb, amsthm, hyperref,
% and a theorem environment. No labels from either manuscript are required.
\paragraph{Stochastic retention with a controlled energy budget.}
The argument combines a standard variable-metric descent estimate
(\href{https://doi.org/10.1137/15M1053141}{Wang et al., 2017, Section 2})
with the nonnegative-supermartingale maximal inequality
(\href{https://doi.org/10.1214/18-PS321}{Howard et al., 2020, Lemma 1}).
The probability consequence for protected exemplars is derived here; it is
not an existing theorem about Adam or language-model replay.

\begin{theorem}[Optimizer-aware conditional exemplar retention]
Fix finitely many complete verified responses $(x_j,e_j)$ with keys $b_j$,
lengths $L_j>0$, and weights $\alpha_j>0$. Let
\[
 R(\theta)=\sum_j\frac{\alpha_j}{L_j}
             [-\log\pi_\theta(e_j\mid x_j)],\qquad
 F(\theta)=R(\theta)-J(\theta),\qquad J(\theta)\le J_{\max}<\infty.
\]
Assume $F$ is differentiable with globally $L_F$-Lipschitz gradient,
and take deterministic $\theta_0$ with finite $F(\theta_0)$. Relative to a filtration
$(\mathcal F_t)$, consider
\[
 \theta_{t+1}=\theta_t-\eta_t H_t(g_t+b_t+\xi_t),\qquad
 g_t=\nabla F(\theta_t).
\]
Here $\theta_t,\eta_t,H_t,b_t$ are $\mathcal F_t$-measurable,
$0\preceq H_t\preceq MI$, $\eta_t>0$, $\eta_tL_FM\le1$, and
$\mathbb E[\xi_t\mid\mathcal F_t]=0$ with finite conditional second moment.
Suppose deterministic numbers $d_t\ge0$ bound
\[
 \frac{\eta_t}{2}\|b_t\|_{H_t}^2
 +\frac{L_F\eta_t^2}{2}
       \mathbb E[\|H_t\xi_t\|_2^2\mid\mathcal F_t]\le d_t,
 \qquad \|u\|_H^2=u^\top Hu.
\]
For $N\in\mathbb N\cup\{\infty\}$ with
$B_N=\sum_{t<N}d_t<\infty$, put
$D_0=F(\theta_0)+J_{\max}\ge0$. For every $0<\delta<1$,
\[
 \Pr\!\left(
  \forall t\le N,\ \forall j:\quad
  p_{\theta_t}(b_j\mid x_j)
  \ge\pi_{\theta_t}(e_j\mid x_j)
  \ge\exp\!\left[-\frac{L_j(D_0+B_N)}{\alpha_j\delta}\right]
 \right)\ge1-\delta.
\]
When $N=\infty$, each protected probability also has an almost surely
positive, potentially random, infimum over all iterations.
\end{theorem}

\begin{proof}
Smoothness, conditional centering of $\xi_t$, and
$H_t^2\preceq MH_t$ give
\begin{align*}
 \mathbb E[F(\theta_{t+1})\mid\mathcal F_t]
 &\le F(\theta_t)-\eta_tg_t^\top H_t(g_t+b_t)
       +\frac{\eta_t}{2}\|g_t+b_t\|_{H_t}^2
       +\frac{L_F\eta_t^2}{2}
          \mathbb E[\|H_t\xi_t\|_2^2\mid\mathcal F_t]\\
 &\le F(\theta_t)-\frac{\eta_t}{2}\|g_t\|_{H_t}^2+d_t.
\end{align*}
The cross term involving $b_t$ cancels on expanding the square.
Since $F+J_{\max}\ge R\ge0$, the adapted process
\[
 X_t=F(\theta_t)+J_{\max}+\sum_{s=t}^{N-1}d_s
\]
is a nonnegative supermartingale; for $N=\infty$ use the infinite tail.
Its initial value is $D_0+B_N$. The maximal inequality yields
$\Pr(\sup_{t\le N}X_t\ge(D_0+B_N)/\delta)\le\delta$.
If $D_0+B_N=0$, nonnegativity instead gives $X_t=0$ almost surely.
Every nonnegative summand of $R(\theta_t)\le X_t$ obeys the same bound,
and exponentiation proves the simultaneous exemplar floors without a
union factor over $j$. Each complete-exemplar event is contained in its
verified-key event. For $N=\infty$, the same maximal inequality, with a
threshold tending to infinity, gives $\sup_tX_t<\infty$ almost surely.
\end{proof}

\paragraph{Finite and infinite horizons are different.}
Unbiased bounded-variance updates have $d_t=O(\eta_t^2)$; square-summable
steps therefore suffice for the infinite-horizon energy budget.
A persistent bias instead costs $O(\eta_t\|b_t\|^2)$.
Constant steps and persistent variance give only the stated finite-horizon
bound unless additional structure controls accumulated noise.
No positive lower eigenvalue of $H_t$ is needed for retention alone;
this theorem does not assert parameter convergence or optimality.
The likelihood must include the complete response and termination event
under the sampling law whose key probability is bounded.

\paragraph{Scope for adaptive implementations.}
A current-gradient Adam second moment is not predictable with respect to
that gradient's noise, and momentum, clipping, bank changes and cyclic
prompt scheduling need separate error or energy control. Positive replay
weight and a small nominal learning rate do not verify these assumptions.
\href{https://arxiv.org/html/1904.09237v1}{Reddi et al. (2018)} provide
counterexamples to unrestricted Adam convergence; this does not itself
prove replay failure in the present application. The bound above is
conditional and can be numerically weak, particularly for long exemplars
or small replay weights.
